\section{Discussion}

\subsection{Why Stage~3 and Stage~7 Are Highest-Leverage}

The asymmetry of failure costs is the central design insight of the pipeline. A puzzle concept that is rejected at Stage~3 (pool ranking) costs only the work of writing a brief — typically a few minutes of prompting and review. A puzzle that clears Stage~3 but fails at or after Stage~7 (editorial review) wastes the full authoring cycle: puzzle construction, answer confirmation, extraction mechanism design, and any associated world-building, amounting to 30–60 minutes of invested effort per puzzle. Two review gates rather than one is therefore not an aesthetic choice but an economic one: the first gate is cheap insurance against expensive downstream work; the second gate is the last high-signal checkpoint before integration costs compound further.

The empirical data is consistent with this framing. Across the 52 puzzles in the five core scenarios, zero failures were detected at Stage~3. This does not mean Stage~3 catches nothing — it means the pool-ranking pass correctly elevated strong candidates and suppressed weak ones before any authoring investment was made. The three constructive failures (Age of Empires Puzzle~V, Ironfall~P01, Ironfall~P02) emerged at Stage~8 (testing and integration), after Stage~7 editorial review had passed them. These failures are attributable not to editorial blindness but to failure modes that only become visible during answer-path verification: a pure-computation mechanism with no deductive layer, and two answer-coordination mismatches arising from brief-to-meta handoff errors. Both are classes of failure that Stage~7 rubrics can be extended to detect, and both are recoverable at Stage~8 with moderate cost. The remaining two failures were a construction error caught at Stage~6 and a delivery-format gap at Stage~9 — neither attributable to editorial gate failure.

The key operational recommendation from this analysis is direct: do not commit to authoring until Stage~3 ranking is complete, and do not commit to integration until Stage~7 editorial verdict is in hand. At any production scale, the cost savings per filtered puzzle are material.

\subsection{The World Lock as a Coherence Mechanism}

Stage~2 produces a canonical world document — \texttt{LOCKED.md} in the pipeline's artifact convention — which is treated as read-only by all subsequent stages. This single-document, single-commit mechanism is the pipeline's primary defense against generation drift: the gradual incoherence that accumulates when later-stage content is generated against a world that has itself been revised.

Without a world lock, the failure mode is subtle and insidious. A character name changed in a late Stage~6 revision would invalidate puzzle flavor text written in Stage~4; a geography change would break a Stage~5 meta mechanism that depended on location names as answer components. These failures are difficult to detect in isolation because each individual artifact may be internally consistent, while the collection as a whole is not. In a multi-stage generative pipeline operating across parallel authoring sessions — as is the case in Stages~5 and~6 of this pipeline — incoherence accumulates silently across agent contexts.

None of the five core scenarios experienced coherence failures attributable to world drift. This negative result is, we argue, precisely the point: the lock worked as intended and the failure class it prevents did not occur. The cost of the lock is negligible — it requires nothing beyond what a scope and structure stage would produce anyway, formalized as a committed artifact. The return is immunity to an entire failure class.

For practitioners adapting the pipeline to new domains, the implication is clear: produce the world lock before running any generative stage, commit it to version control, and communicate its read-only status explicitly in the scenario's \texttt{CLAUDE.md}. The one-time setup cost is dominated by the coherence insurance it provides.

\subsection{The Autonomy Gradient}

A counterintuitive finding from the autonomy coding analysis (Section~4) is that the pipeline's middle stages — authoring (Stage~6) and editorial review (Stage~7) — are where AI operates most autonomously, while the earliest and latest stages involve the greatest proportion of human judgment. Stage~6 in particular was executed entirely by AI agents across all scenarios, with parallel agent launches and no human editing of puzzle files during construction. Stage~7 editorial review was AI-panel-driven in all cases, with human involvement limited to disposition decisions on borderline verdicts.

This gradient is counterintuitive because creative work is typically assumed to require human oversight precisely at the point of creation. The explanation is structural: the middle stages are where domain knowledge is densest and most tractable for AI systems. Stage~6 operates against a fully specified brief (from Stage~4), a locked world (from Stage~2), and a known answer (from Stage~5). The generation task is highly constrained, and the quality rubric is explicit. Stage~7 applies a formalized editorial rubric against a completed artifact — a task well-suited to multi-profile panel evaluation. Constraint density, not creativity, is what drives AI autonomy.

By contrast, Stage~1 (scope definition) requires human judgment about intent — what the hunt should \emph{mean}, what experience it should produce, who the audience is. Stage~11 (polish) requires judgment about edge cases the system did not anticipate: hint calibration for specific solver populations, accessibility accommodations, last-mile delivery decisions. Neither stage admits the kind of structured task decomposition that makes AI autonomous execution reliable. Human involvement is most valuable at the boundaries where structure is thinnest.

This gradient has a practical implication for pipeline adoption: practitioners should expect to invest human time at Stage~1 and Stage~11, and should trust AI autonomy most heavily in the middle. Attempts to increase human involvement in Stage~6 authoring, perhaps out of skepticism about AI-generated puzzle content, are likely to add cost without proportionate quality benefit.

\subsection{Limitations}

Several limitations constrain the generalizability of these findings.

\textbf{Single toolkit system.} All five scenarios were executed using the same \texttt{CLAUDE.md}-driven skill library and reviewer profile set. Cross-pipeline transferability — whether a different authoring system with different underlying LLM prompts would yield comparable stage-level pass rates — is not demonstrated here. The pipeline \emph{specification} is transferable; the empirical results are specific to this implementation.

\textbf{AI panel as quality proxy.} Pass rates throughout this paper are assessed by AI expert panel review, not by human play-testing. Whether AI panel verdict at Stage~7 correlates with solver enjoyment, puzzle fairness, or aha-moment quality is studied in a companion paper~\cite{dellaLibera2025calibration}. The results reported here establish pipeline behavior under the AI-panel quality criterion; they should not be taken as evidence about human solver experience without that validation. All quality metrics are derived from AI panel review; correlation with solver enjoyment, engagement, and perceived difficulty is studied in companion work~\cite{dellaLibera2025calibration}. The pipeline optimizes for panel-pass quality, which is a proxy for but not identical to player experience.

\textbf{Maturation confound.} The 90.4\% aggregate pass rate spans a materially evolving pipeline; improvement attributable to pipeline refinement, larger candidate pools, and accumulated session experience cannot be disentangled.

\textbf{Absence of baseline condition.} We do not report a gate-disabled baseline condition; the causal attribution of quality improvement to pipeline structure is inferential rather than experimentally established.

\textbf{Scale ceiling.} The largest scenario in the empirical set is Dead Reckoning at 19 puzzles (15 feeders, 4 metas). MIT Mystery Hunt scale — 200 or more puzzles with nested meta structure and cross-round interactions — is entirely untested. Pipeline behavior at that scale, particularly the coherence properties of the world lock and the compounding effects of Stage~8 answer-coordination failures, is an open empirical question.

\textbf{Timing as proxy.} Stage timing data is derived from commit timestamps, which capture wall-clock intervals between artifact commits rather than active human-minutes per stage. Stages executed in overlapping sessions (Scenarios~3 through~6 ran in parallel contexts on Day~1) compound this measurement limitation. The timing analysis presented in Section~4 should be understood as an order-of-magnitude characterization, not a precise human-effort accounting.

\subsection{Design Recommendations for Practitioners}

Five recommendations follow from the empirical analysis.

\textbf{Run Stage~3 pool ranking before committing to authoring.} The cost savings per filtered puzzle are significant at any scale. A pool of 15 candidates ranked by an AI panel requires less than 20 minutes (as observed in Stage~3 of the Age of Empires scenario at 8–20 minutes elapsed); committing to authoring without ranking exposes all candidate slots to the full authoring cost regardless of quality.

\textbf{Invest in the world-lock document before running generative stages.} The lock prevents coherence failures at no additional cost beyond what the scope and structure stages produce in the normal course. The one-time overhead of committing \texttt{LOCKED.md} to version control is dominated by the coherence insurance it provides across all subsequent stages.

\textbf{Run AI expert panels on every puzzle candidate, not only selected finalists.} Panel evaluation is cheap enough — and reliable enough relative to the cost of a wasted authoring cycle — that broad application is warranted. Reserving panel review for a pre-selected shortlist defeats the purpose of the Stage~3 gate.

\textbf{Configure per-scenario scope and world documents; reuse the skill library unchanged.} The five-scenario empirical record confirms that the \texttt{CLAUDE.md} skill library transfers without modification across domains, scales, and delivery formats. Per-scenario effort should be concentrated in Stage~1 (scope) and Stage~2 (world design), not in pipeline reconfiguration.

\textbf{Expect and plan for Stage~8 failures.} Three of the five pipeline failures surfaced at Stage~8, after editorial review had passed the puzzles. A small Stage~8 failure budget (10\% of final puzzle count) should be assumed when estimating delivery timelines, and the pipeline should be operated with sufficient puzzle-pool depth to absorb redesigns without missing the final count target.
